\documentclass[11pt,a4paper]{article}
\usepackage[utf8]{inputenc}
\usepackage[T1]{fontenc}
\usepackage{lmodern}
\usepackage[french]{babel}
\usepackage{amsmath,amssymb,mathtools}
\usepackage{icomma}
\usepackage{xcolor}
\usepackage{tikz}
\usepackage{pgfplots}
\usepackage[most]{tcolorbox}
\usepackage{enumitem}
\usepackage{fancyhdr}
\usepackage[hidelinks]{hyperref}
\usepackage{titlesec}
\usepackage[a4paper,margin=2.2cm,headheight=26pt,headsep=10pt]{geometry}
\usetikzlibrary{arrows.meta,positioning,calc,intersections,angles,quotes}
\usepgfplotslibrary{fillbetween}
\pgfplotsset{compat=1.18}
\definecolor{primary}{RGB}{30,75,145}
\definecolor{primarydark}{RGB}{20,50,100}
\definecolor{defcol}{RGB}{0,114,178}
\definecolor{thmcol}{RGB}{0,135,110}
\definecolor{excol}{RGB}{0,130,85}
\definecolor{remarkcol}{RGB}{120,70,180}
\definecolor{attncol}{RGB}{200,40,55}
\definecolor{methcol}{RGB}{85,85,100}
\definecolor{areacol}{RGB}{120,160,230}
\definecolor{areacol2}{RGB}{230,150,80}
\newcommand{\dd}{\mathrm{d}}
\newcommand{\R}{\mathbb{R}}
\newcommand{\ua}{\,\text{u.a.}}
\newcommand{\tg}{\tan}\newcommand{\cotg}{\cot}
\tcbset{boxbase/.style={enhanced,breakable,boxrule=0.9pt,arc=2.5pt,left=3mm,right=3mm,top=2.3mm,bottom=2.3mm,fonttitle=\bfseries,coltitle=white,toptitle=0.6mm,bottomtitle=0.6mm}}
\newtcolorbox{methode}[1][]{boxbase,colback=methcol!7,colframe=methcol,sharp corners,title={\textsc{Comment faire}\ifx\relax#1\relax\relax\else\textmd{ : \itshape #1}\fi}}
\newtcolorbox{exemple}[1][]{boxbase,colback=excol!5,colframe=excol,title={\textsc{Exemple}\ifx\relax#1\relax\relax\else\textmd{ : \itshape #1}\fi}}
\newtcolorbox{idee}[1][]{boxbase,colback=defcol!5,colframe=defcol,title={\textsc{Idée clé}\ifx\relax#1\relax\relax\else\textmd{ : \itshape #1}\fi}}
\newtcolorbox{remarque}[1][]{boxbase,colback=remarkcol!6,colframe=remarkcol,title={\textsc{Remarque}\ifx\relax#1\relax\relax\else\textmd{ : \itshape #1}\fi}}
\newtcolorbox{attention}[1][]{boxbase,colback=attncol!6,colframe=attncol,title={\textsc{Attention}\ifx\relax#1\relax\relax\else\textmd{ : \itshape #1}\fi}}
\newtcolorbox{exo}[1][]{boxbase,colback={primary!4},colframe=primary,title={\textsc{Exercice #1}}}
\newtcolorbox{corr}[1][]{boxbase,colback={thmcol!5},colframe=thmcol,title={\textsc{Corrigé #1}}}
\tikzset{axe/.style={->,>=Stealth,black!65,line width=0.7pt},courbe/.style={primary,line width=1.3pt},courbe2/.style={areacol2!90!black,line width=1.3pt},dvert/.style={gray,dashed,line width=0.7pt}}
\pagestyle{fancy}\fancyhf{}
\fancyhead[L]{\small\textcolor{primarydark}{\textbf{Fiche 7} — Intégrale d'une fonction continue}}
\fancyhead[R]{\small\textcolor{primarydark}{\textsc{Thème 3 — Fractions rationnelles}}}
\fancyfoot[C]{\small\thepage}
\renewcommand{\headrulewidth}{0.8pt}
\titleformat{\section}{\large\bfseries\color{primarydark}}{\thesection}{0.6em}{}
\begin{document}
\section*{Thème 3 — Fractions rationnelles \quad\textbullet\quad Corrigé Moyen}

\begin{corr}[1]
$\dfrac{2x}{(x+1)(x+3)}=\dfrac{A}{x+1}+\dfrac{B}{x+3}$.
$A=\left.\dfrac{2x}{x+3}\right|_{x=-1}=\dfrac{-2}{2}=-1$,
$B=\left.\dfrac{2x}{x+1}\right|_{x=-3}=\dfrac{-6}{-2}=3$.
\emph{Contrôle :} $-1(x+3)+3(x+1)=-x-3+3x+3=2x$ \checkmark. Donc
\[
\int \frac{2x}{(x+1)(x+3)}\,\dd x=-\ln|x+1|+3\ln|x+3|+C.
\]
\end{corr}

\begin{corr}[2]
$\dfrac{2x-1}{(x+2)(x-3)}=\dfrac{A}{x+2}+\dfrac{B}{x-3}$.
$A=\left.\dfrac{2x-1}{x-3}\right|_{x=-2}=\dfrac{-5}{-5}=1$,
$B=\left.\dfrac{2x-1}{x+2}\right|_{x=3}=\dfrac{5}{5}=1$.
\emph{Contrôle :} $1(x-3)+1(x+2)=2x-1$ \checkmark. Donc
\[
\int \frac{2x-1}{(x+2)(x-3)}\,\dd x=\ln|x+2|+\ln|x-3|+C.
\]
\end{corr}

\begin{corr}[3]
\textbf{Factorisation.} $x^2-x-6$ a pour racines $3$ et $-2$ (somme $1$, produit $-6$), donc
$x^2-x-6=(x-3)(x+2)$.

\smallskip
$\dfrac{x+7}{(x-3)(x+2)}=\dfrac{A}{x-3}+\dfrac{B}{x+2}$.
$A=\left.\dfrac{x+7}{x+2}\right|_{x=3}=\dfrac{10}{5}=2$,
$B=\left.\dfrac{x+7}{x-3}\right|_{x=-2}=\dfrac{5}{-5}=-1$.
\emph{Contrôle :} $2(x+2)-1(x-3)=2x+4-x+3=x+7$ \checkmark. Donc
\[
\int \frac{x+7}{x^2-x-6}\,\dd x=2\ln|x-3|-\ln|x+2|+C.
\]
\end{corr}

\begin{corr}[4]
\textbf{Factorisation.} $x^2-4=(x-2)(x+2)$.

\smallskip
$\dfrac{x}{(x-2)(x+2)}=\dfrac{A}{x-2}+\dfrac{B}{x+2}$.
$A=\left.\dfrac{x}{x+2}\right|_{x=2}=\dfrac{2}{4}=\tfrac12$,
$B=\left.\dfrac{x}{x-2}\right|_{x=-2}=\dfrac{-2}{-4}=\tfrac12$.
\emph{Contrôle :} $\tfrac12(x+2)+\tfrac12(x-2)=x$ \checkmark. Donc
\[
\int \frac{x}{x^2-4}\,\dd x=\frac12\ln|x-2|+\frac12\ln|x+2|+C
=\frac12\ln|x^2-4|+C.
\]
\end{corr}

\begin{corr}[5]
\textbf{Factorisation.} $x^2+x-2$ a pour racines $1$ et $-2$ (somme $-1$, produit $-2$), donc
$x^2+x-2=(x-1)(x+2)$.

\smallskip
$\dfrac{3x+1}{(x-1)(x+2)}=\dfrac{A}{x-1}+\dfrac{B}{x+2}$.
$A=\left.\dfrac{3x+1}{x+2}\right|_{x=1}=\dfrac{4}{3}$,
$B=\left.\dfrac{3x+1}{x-1}\right|_{x=-2}=\dfrac{-5}{-3}=\dfrac53$.
\emph{Contrôle :} $\tfrac43(x+2)+\tfrac53(x-1)=\tfrac43x+\tfrac83+\tfrac53x-\tfrac53
=\tfrac{9}{3}x+\tfrac{3}{3}=3x+1$ \checkmark. Donc
\[
\int \frac{3x+1}{x^2+x-2}\,\dd x=\frac43\ln|x-1|+\frac53\ln|x+2|+C.
\]
\end{corr}

\end{document}
